%%%%%%%%%%%%%%%%%%%%%%%%%%%%%%%%%%%%%%%%%%%%%%%%%%%%%%%%%%%%%%%%%%%%%%%%%%%%%%%%%%%%%%%%%%%%%%%%%%%%%%%%%%%%%%%
%%%%%%%%%%%%%%%%%%%%%%%%%%%%%%%%%%%%%%%%%%%%%%%%%%%%%%%%%%%%%%%%%%%%%%%%%%%%%%%%%%%%%%%%%%%%%%%%%%%%%%%%%%%%%%%
% Esta é apenas uma adaptação do ABNTex-2 para ser utilizado como tese ou dissertação de mestrado em inglês.
% Em resumo, os termos foram adptados para inglês (ex: Orietador se tornou Supervisor) e o modelo de citação, ao invés de 
% seguir um padrão todo maiusculo, foi trocado para minusculo (ex: FULANO; CICRANO; BELTRANO, 2010 se tornou Fulano, Cicrano, and Beltrano, % % 2010.
% Você pode se perguntar: Pra que usar o formato ABNT e inglês? Bom, algumas universidades não tem um padrão definido de tese/dissertação,
% e o ABNTex, com esses ajustes, fica um padrão bacana para ser seguido. Essas alterações visam poupar o seu tempo para configurar um projeto.
% Por fim, todos os créditos para a equipe que desenvolveu o pacote. Isso é apenas uma adaptação minuscula.
% Caso você venha a utilizar essa adptação, considere deixar a sua estrelinha para o projeto no Github, 
% isso valoriza e motiva desenvolvedores.
% Github: https://github.com/abntex/abntex2
% Faça bom uso!
% André Pacheco (pacheco.comp@gmail.com)
%%%%%%%%%%%%%%%%%%%%%%%%%%%%%%%%%%%%%%%%%%%%%%%%%%%%%%%%%%%%%%%%%%%%%%%%%%%%%%%%%%%%%%%%%%%%%%%%%%%%%%%%%%%%%%%
%%%%%%%%%%%%%%%%%%%%%%%%%%%%%%%%%%%%%%%%%%%%%%%%%%%%%%%%%%%%%%%%%%%%%%%%%%%%%%%%%%%%%%%%%%%%%%%%%%%%%%%%%%%%%%%

\documentclass[12pt, openright, oneside, a4paper, english]{abntex2}



% ---
% Pacotes básicos 
% ---
\usepackage{lmodern}			% Usa a fonte Latin Modern			
\usepackage[T1]{fontenc}		% Selecao de codigos de fonte.
\usepackage[utf8]{inputenc}		% Codificacao do documento (conversão automática dos acentos)
\usepackage{lastpage}			% Usado pela Ficha catalográfica
\usepackage{indentfirst}		% Indenta o primeiro parágrafo de cada seção.
\usepackage{color}				% Controle das cores
\usepackage{graphicx}			% Inclusão de gráficos
\usepackage{microtype} 			% para melhorias de justificação
\usepackage{amsthm}
\usepackage{amsmath}
\usepackage{mathtools}
\usepackage{amsfonts}
\usepackage{amssymb}
%\usepackage{mathtools}
\usepackage{caption}
\usepackage{subcaption}
\usepackage{tabularx}
\usepackage{csquotes}
\usepackage{makecell}

% \usepackage{subfigure}

%\usepackage{kbordermatrix}
\usepackage{amsmath}

% \usepackage[portuguese,ruled,lined,linesnumbered]{algorithm2e}
\usepackage[ruled,lined,linesnumbered]{algorithm2e}
% \renewcommand{\algorithmcfname}{Algorithm}

\usepackage{multirow}

\usepackage[brazil]{babel}
\addto\captionsbrazil{%
\def\bibname{References}%
\renewcommand{\figurename}{Figure}
\renewcommand{\tablename}{Table}
\renewcommand{\listfigurename}{List of Figures}
\renewcommand{\listtablename}{List of Tables}
\renewcommand{\contentsname}{Contents}
\renewcommand\chaptername{Chapter}
}

\selectlanguage{english}


% ---
		
% ---
% Pacotes adicionais, usados apenas no âmbito do Modelo Canônico do abnteX2
% ---
\usepackage{lipsum}				% para geração de dummy text
% ---

% ---
% Pacotes de citações
% ---
%\usepackage{backref}	 % Paginas com as citações na bibl
%\usepackage[alf,abnt-and-type=&]{abntex2cite}	% Citações padrão ABNT
\usepackage[alf]{abntex2cite}	% Citações padrão ABNT


%\newtheorem{mydef}{Definição}[chapter]

%\usepackage{natbib}

% --- 
% CONFIGURAÇÕES DE PACOTES
% --- 

% ---
% Configurações do pacote backref
% Usado sem a opção hyperpageref de backref
%\renewcommand{\backrefpagesname}{Citado na(s) página(s):~}
% Texto padrão antes do número das páginas
%\renewcommand{\backref}{}
% Define os textos da citação
%\renewcommand*{\backrefalt}[4]{
%	\ifcase #1 %
%		Nenhuma citação no texto.%
%	\or
%		Citado na página #2.%
%	\else
%		Citado #1 vezes nas páginas #2.%
%	\fi}%
% ---
\newtheorem{name}{Printed output}

\selectlanguage{english}
\renewcommand{\orientadorname}{Supervisor}







% ---
% Informações de dados para CAPA e FOLHA DE ROSTO
% ---
\titulo{On the power of invisibility using a simple ring}
\autor{Bilbo Baggins}
\local{The Shire}
\data{2020}
\orientador{Gandalf, the White}

\instituicao{%
 University of the Shire
  \par
 Graduate Program in Old Sorcery }
\tipotrabalho{Doctorate Thesis}
% O preambulo deve conter o tipo do trabalho, o objetivo, 
% o nome da instituição e a área de concentração 
\preambulo{Doctoral Thesis presented to the Graduate Program in Old Sorcery of the University of the Shire as a partial requirement to obtain the degree of Doctor in Old Sorcery}
% ---


% ---
% Configurações de aparência do PDF final

% alterando o aspecto da cor azul
\definecolor{blue}{RGB}{41,5,195}

% informações do PDF
\makeatletter
\hypersetup{
     	%pagebackref=true,
		pdftitle={\@title}, 
		pdfauthor={\@author},
    	pdfsubject={\imprimirpreambulo},
	    pdfcreator={LaTeX with abnTeX2},
		pdfkeywords={abnt}{latex}{abntex}{abntex2}{trabalho acadêmico}, 
		colorlinks=true,       		% false: boxed links; true: colored links
    	linkcolor=blue,          	% color of internal links
    	citecolor=blue,        		% color of links to bibliography
    	filecolor=magenta,      		% color of file links
		urlcolor=blue,
		bookmarksdepth=4
}
\makeatother

% O tamanho do parágrafo é dado por:
\setlength{\parindent}{1.3cm}

% Controle do espaçamento entre um parágrafo e outro:
\setlength{\parskip}{0.2cm}  % tente também \onelineskip

% ---
% compila o indice
% ---
\makeindex

% ----
\begin{document}

% Retira espaço extra obsoleto entre as frases.
\frenchspacing 

% ----------------------------------------------------------
% ELEMENTOS PRÉ-TEXTUAIS
% ----------------------------------------------------------
\pretextual

% ---
% Capa
% ---
\imprimircapa
% ---
\cleardoublepage

% ---
% Folha de rosto
% (o * indica que haverá a ficha bibliográfica)
% ---
\imprimirfolhaderosto
% ---

% ---
% Inserir a ficha bibliografica
% ---
% \begin{fichacatalografica}
% 	\begin{figure}
% 	 \includegraphics[scale=0.9]{ficha2.pdf}
% 	\end{figure}
% \end{fichacatalografica}


% \begin{folhadeaprovacao}
%     \begin{figure}
%      \includegraphics[bb=110 0 0 750]{aprov.pdf}
%     \end{figure}
% \end{folhadeaprovacao}


% ---
% Declaração de autoria
% ---
\newpage
\begin{center}
{\ABNTEXchapterfont\Large\textsc{Statement of Authorship}}
\end{center}

\vspace*{1cm}

I hereby declare that the thesis submitted is my own work. All direct or indirect
sources used are acknowledged as references. I further declare that I have not submitted
this thesis at any other institution in order to obtain a degree. \\

% ---


\newpage
% ---
% Dedicatória
% ---
\begin{dedicatoria}
  \vspace*{\fill}
  \centering
  \noindent
  \textit{ To my friends and family. I am nothing without them.} \vspace*{\fill}
\end{dedicatoria}
% ---



% ---
% Agradecimentos
% ---
\begin{agradecimentos}

I would like to thank all people from The Shire, in particular, my beloved nephew Frodo.

\end{agradecimentos}
% ---

\newpage
% ---
% Epígrafe
% ---
\begin{epigrafe}
    \vspace*{\fill}
	
	    \noindent
		\textit{"One Ring to rule them all, One Ring to find them, One Ring to bring them all and in the darkness bind them.."}
		
		\begin{flushright}
		J. R. R. Tolkien
		\end{flushright}
\end{epigrafe}


% ---
\newpage

 

 

% resumo em português
\setlength{\absparsep}{18pt} % ajusta o espaçamento dos parágrafos do resumo
\begin{resumo}[Abstract]
 \begin{otherlanguage*}{english}


\input{src/abstract}


\textbf{Keywords:} Invisibility, Sorcery, Ring

\end{otherlanguage*}
\end{resumo}

\newpage


% ---
% inserir lista de ilustrações
% ---
\pdfbookmark[0]{\listfigurename}{lof}
\listoffigures*
\cleardoublepage
% ---

% ---
% inserir lista de tabelas
% ---
\pdfbookmark[0]{\listtablename}{lot}
\listoftables*
\cleardoublepage
% ---
%\begin{siglas}
  %\item[MSE] Erro médio quadrático
  %\item[FLR] Relacionamento Lógico Fuzzy
%\end{siglas}

% ---
% inserir o sumario
% ---
\pdfbookmark[0]{\contentsname}{toc}
\tableofcontents*
\cleardoublepage
% ---

% ----------------------------------------------------------
% ELEMENTOS TEXTUAIS
% ----------------------------------------------------------
\textual

% ----------------------------------------------------------
% Introdução (exemplo de capítulo sem numeração, mas presente no Sumário)
% ----------------------------------------------------------
\chapter{Introduction}
\input{src/intro}

\chapter{Raug}
\input{src/another_chapter}



% ----------------------------------------------------------
% Referências bibliográficas
% ----------------------------------------------------------
\citeoption{abnt-etal-list=5}
\bibliographystyle{abntex2-cite-min}
\bibliography{references}



% ----------------------------------------------------------
% Glossário
% ----------------------------------------------------------
%
% Consulte o manual da classe abntex2 para orientações sobre o glossário.
%
%\glossary

% ----------------------------------------------------------
% Apêndices
% ----------------------------------------------------------

% ---
% Inicia os apêndices
% ---
\begin{apendicesenv}

% Imprime uma página indicando o início dos apêndices
%\partapendices

% ----------------------------------------------------------
\chapter{People who have worn the ring}
% ----------------------------------------------------------
\input{src/appendice}

% ----------------------------------------------------------
%\chapter{Nullam elementum urna vel imperdiet sodales elit ipsum pharetra ligula
%ac pretium ante justo a nulla curabitur tristique arcu eu metus}
% ----------------------------------------------------------
%\lipsum[55-57]

\end{apendicesenv}
% ---


% ----------------------------------------------------------
% Anexos
% ----------------------------------------------------------

% ---
% Inicia os anexos
% ---
%\begin{anexosenv}

% Imprime uma página indicando o início dos anexos
%\partanexos


%\end{anexosenv}

%---------------------------------------------------------------------
% INDICE REMISSIVO
%---------------------------------------------------------------------
\phantompart
\printindex
%---------------------------------------------------------------------

\end{document}
